\documentclass[11pt,a4paper]{article}

\usepackage[margin=2.4cm]{geometry}
\usepackage{graphicx}
\usepackage{booktabs}
\usepackage{hyperref}
\usepackage{caption}

\graphicspath{{../analysis/out/}}
\hypersetup{colorlinks=true, linkcolor=black, urlcolor=blue, citecolor=black}

\title{\textbf{Split Inference on a Heterogeneous ESP32 Fleet}\\[0.3em]
\large Project Work in Internet of Things}
\author{Giacomo Antonelli}
\date{September 2026}

\begin{document}
\maketitle

\begin{abstract}
\noindent
Split inference runs the first $k$ layers of a neural network on a
microcontroller and sends the intermediate activation tensor to a server that
finishes the job. We built the full pipeline on three ESP32-class boards of
differing capability and measured what the choice of $k$ costs. Across
3{,}813 requests with zero corrupted deliveries, full offload beat every split
point on every board, by $5.1\times$ on an ESP32-S3 and $15.5\times$ on an
ESP32-CAM, and the weakest board could not host a split point at all. Cutting
deeper made things worse rather than better: the deeper cut shrank the tensor
from 4{,}668\,B to 888\,B but bought 2.7\,ms of transfer with 26--164\,ms of
extra on-device inference. The per-stage breakdown shows why. Device compute
dominates by two orders of magnitude, while transfer, the quantity split
inference exists to reduce, costs 7--12\,ms regardless of which cut produced it.
A sensitivity analysis over the measured stages puts break-even at a server
$51\times$ slower than ours.
\end{abstract}

\section{The question}

For a network layer $k$, a split deployment pays three costs: on-device
inference of the first $k$ layers, transfer of the resulting activation tensor,
and server inference of the remainder. Cutting later means more device compute
and usually less data on the air. The premise of split inference is that some
intermediate $k$ beats both endpoints, and that the best $k$ depends on the
device.

This report tests that premise by measurement. It does not learn or adapt the
split point: every policy here is a fixed rule, and the exhaustive sweep is the
ground truth an adaptive orchestrator would have to beat.

\section{System}

\subsection{Hardware}

\begin{table}[h]
\centering
\begin{tabular}{llll}
\toprule
Node & Tier & Board & Memory \\
\midrule
11 & A & ESP32-S3-WROOM-1 N16R8 & 8\,MB octal PSRAM @ 80\,MHz, vector ISA \\
21 & B & ESP32-CAM (ESP32-D0WD)  & 8\,MB quad PSRAM @ 80\,MHz, no vector ISA \\
31 & C & ESP32-D0WD-V3, plain    & no PSRAM, ${\sim}150$\,KB usable SRAM \\
\midrule
   &   & x86 laptop (Core 7 150U) & soft-AP, orchestrator, server tail \\
\bottomrule
\end{tabular}
\caption{The fleet. All three MCUs run at 240\,MHz.}
\end{table}

All boards run ESP-IDF v5.5.3 with TensorFlow Lite Micro. The server is the same
machine that builds the firmware: it runs the Wi-Fi access point (2.4\,GHz, since
ESP32 radios have no 5\,GHz band), the orchestrator, and the tail inference. Its
CPU is capped to 400\,MHz and the process pinned to four cores. Section
\ref{sec:threats} states what that costs the conclusions.

\subsection{Model and split points}

MobileNetV2, width multiplier 0.35, $96\times96\times3$ input, full-int8
post-training quantization, two classes. Heads and tails are sliced from one set
of weights at the flatbuffer level, with the boundary tensor's quantization
parameters copied verbatim from head output to tail input. A test verifies
bit-exact equality of int8 logits between $\mathrm{full}(x)$ and
$\mathrm{tail}_k(\mathrm{head}_k(x))$ over 200 images at every cut, so accuracy
is identical by construction and latency is the only thing that varies between
policies.

Three split points: \texttt{k0} (no split, meaning JPEG to the server and the
whole model on the server), \texttt{k\_shallow} (after inverted-residual block
2), and \texttt{k\_deep} (after block 6).

\subsection{Transport and measurement}

UDP with application-level fragmentation, CRC32 per tensor, NACK-based selective
repeat over three rounds, and an abort path. One asyncio process on the server
dispatches at most one request per node at a time. Every request writes one
SQLite row decomposing end-to-end latency into five non-overlapping stages that
sum exactly to the total: \texttt{device} (flash read plus head inference),
\texttt{uplink}, \texttt{queue}, \texttt{server}, \texttt{residual}. Every number
in this report is a query over those rows.

\section{Method}

Two kinds of run. The \emph{sweep} measures every feasible \emph{(node, cut)}
pair, 200 requests each, one node at a time. That serialisation is what makes
the resulting numbers properties of the cut rather than of what the other nodes
happened to be doing. The \emph{policy runs} then exercise four fixed rules with
the whole fleet active: \texttt{k0}, \texttt{k\_shallow}, \texttt{k\_deep}, and
\texttt{best}, which replays the sweep's per-node argmin.

Every board reads the same 50 images from flash, in the same order, so all
policies see identical inputs.

Seven gates guard the procedure; each invalidates everything after it if it
fails. They cover split correctness, ARQ integrity under injected loss,
whole-harness verification against simulated nodes, per-tier arena feasibility,
fleet stability, sample counts, and dispatch correctness.

\section{Results}

\subsection{Which split points each board can host}

\begin{table}[h]
\centering
\begin{tabular}{llrrl}
\toprule
Tier & Arena source & \texttt{k\_shallow} & \texttt{k\_deep} & Feasible cuts \\
\midrule
A & 8\,MB octal PSRAM   & 149{,}972\,B & 158{,}120\,B & \texttt{k0}, \texttt{k\_shallow}, \texttt{k\_deep} \\
B & 8\,MB quad PSRAM    & 142{,}296\,B & 150{,}396\,B & \texttt{k0}, \texttt{k\_shallow}, \texttt{k\_deep} \\
C & 100\,KB internal    & \multicolumn{2}{c}{needs 138{,}240\,B, has 98{,}480\,B} & \texttt{k0} only \\
\bottomrule
\end{tabular}
\caption{Arena required per head, measured on device at boot. Tier C is short by
39\,KB and hosts no head at all: its firmware links none.}
\label{tab:feasibility}
\end{table}

Two results here. First, the weakest board cannot split, so on that hardware the
only available policy is full offload. Second, the same cut costs different
arena on different silicon: 149{,}972\,B on the S3 against 142{,}296\,B on the
ESP32 for identical weights and tensor shapes, because the S3 build uses vector
kernels that allocate their own scratch.

\subsection{Bytes on the air}

\begin{table}[h]
\centering
\begin{tabular}{lrrrr}
\toprule
Cut & Activation bytes & Mean bytes on air & Fragments & Retransmits/req \\
\midrule
\texttt{k0}        & 1{,}775 (JPEG) & 1{,}811 & 1.98 & 0.000 \\
\texttt{k\_shallow} & 4{,}608 & 4{,}668 & 4.00 & 0.000 \\
\texttt{k\_deep}    & 864     & 888     & 1.00 & 0.000 \\
\bottomrule
\end{tabular}
\caption{The data-inflation effect. Cutting shallow puts $2.6\times$ more on the
air than simply sending the JPEG; only the deep cut falls below it.}
\label{tab:bytes}
\end{table}

\subsection{Latency per node and split point}

\begin{table}[h]
\centering
\begin{tabular}{llrrrr}
\toprule
Node & Cut & $n$ & Mean (ms) & p50 (ms) & p95 (ms) \\
\midrule
11 (A) & \textbf{\texttt{k0}} & 200 & \textbf{32.6} & 29.8 & 55.5 \\
11 (A) & \texttt{k\_shallow}  & 200 & 166.8 & 167.5 & 184.2 \\
11 (A) & \texttt{k\_deep}     & 200 & 190.0 & 189.9 & 209.6 \\
\midrule
21 (B) & \textbf{\texttt{k0}} & 200 & \textbf{36.0} & 33.0 & 63.8 \\
21 (B) & \texttt{k\_shallow}  & 200 & 557.8 & 557.7 & 581.4 \\
21 (B) & \texttt{k\_deep}     & 200 & 719.9 & 720.3 & 743.6 \\
\midrule
31 (C) & \textbf{\texttt{k0}} & 200 & \textbf{41.5} & 40.4 & 68.4 \\
\bottomrule
\end{tabular}
\caption{The sweep: every feasible pair, measured in isolation. Zero failures.}
\label{tab:sweep}
\end{table}

Full offload wins on every node, at the mean and at p95, by $5.1\times$ on tier
A and $15.5\times$ on tier B. And \texttt{k\_deep} is worse than
\texttt{k\_shallow} on both capable tiers, despite putting a fifth as much data
on the air.

\begin{figure}[h]
\centering
\includegraphics[width=0.86\textwidth]{latency_by_cut.png}
\caption{Latency per split point; bar is the mean, rule is p95.}
\end{figure}

\subsection{Where the time goes}

\begin{table}[h]
\centering
\begin{tabular}{llrrrrr}
\toprule
Node & Cut & device & uplink & queue & server & residual \\
\midrule
11 (A) & \texttt{k0}         & 19.4 & 7.6 & 0.1 & 4.3 & 1.1 \\
11 (A) & \texttt{k\_shallow} & \textbf{154.5} & 9.7 & 0.1 & 1.5 & 1.1 \\
11 (A) & \texttt{k\_deep}    & \textbf{180.6} & 7.0 & 0.1 & 1.1 & 1.2 \\
21 (B) & \texttt{k0}         & 22.3 & 8.5 & 0.1 & 4.1 & 1.1 \\
21 (B) & \texttt{k\_shallow} & \textbf{543.9} & 11.0 & 0.1 & 1.6 & 1.2 \\
21 (B) & \texttt{k\_deep}    & \textbf{708.1} & 9.1 & 0.1 & 1.3 & 1.3 \\
31 (C) & \texttt{k0}         & 24.3 & 11.9 & 0.1 & 4.2 & 1.1 \\
\bottomrule
\end{tabular}
\caption{Mean stage breakdown in ms. The stages sum to the total.}
\label{tab:stages}
\end{table}

The \texttt{device} stage dominates every split cut and nothing else is within
two orders of magnitude. Meanwhile \texttt{uplink} is almost constant across
cuts, 7.0 to 11.9\,ms, even though the payload varies $5\times$ between
\texttt{k\_deep} and \texttt{k\_shallow}. At these sizes airtime is per-packet
overhead rather than bytes, so the data-inflation effect of
Table~\ref{tab:bytes} is real and costs almost nothing in time.

\subsection{The cost of running the fleet together}

\begin{table}[h]
\centering
\begin{tabular}{llrrr}
\toprule
Policy & Node & Isolated (ms) & Concurrent (ms) & Difference \\
\midrule
\texttt{k0}          & 11 (A) & 32.6 & 42.4 & \textbf{+30.2\%} \\
\texttt{best} $=$ \texttt{k0} & 11 (A) & 32.6 & 45.8 & \textbf{+40.7\%} \\
\texttt{k\_shallow}  & 11 (A) & 166.8 & 174.1 & +4.4\% \\
\texttt{k\_deep}     & 11 (A) & 190.0 & 193.5 & +1.8\% \\
\bottomrule
\end{tabular}
\caption{The same node running the same cut, measured alone versus with the
whole fleet active. Policies that lean on the shared server tail pay an order of
magnitude more for concurrency than policies that compute on the device.}
\label{tab:coupling}
\end{table}

The three nodes share one single-worker server tail, so a policy whose tail work
is expensive makes every other node wait. The effect is there and in the
predicted direction: splitting genuinely does buy immunity to server contention.
It does not change the ranking, since \texttt{k0} at 42.4\,ms with its penalty
still beats \texttt{k\_shallow} at 174.1\,ms by $4.1\times$.

\subsection{How much slower would the server have to be?}

The server here is far faster than an MCU-class edge device, which favours
offload. Holding the measured \texttt{device}, \texttt{uplink} and
\texttt{residual} stages fixed and scaling the measured \texttt{server} stage
gives the crossover where a split cut would overtake \texttt{k0}:

\begin{table}[h]
\centering
\begin{tabular}{llrr}
\toprule
Node & Takes over & Server slower by & \texttt{k0} tail at crossover \\
\midrule
11 (A) & \texttt{k\_shallow} & $51\times$  & 218\,ms \\
21 (B) & \texttt{k\_shallow} & $217\times$ & 881\,ms \\
31 (C) & none & no crossover & \\
\bottomrule
\end{tabular}
\caption{A projection over measured stages, not a measurement. Queueing is
excluded, which makes every crossover here conservative.}
\label{tab:sensitivity}
\end{table}

For scale, a Raspberry Pi 4, the server this study was originally designed
around, runs this model in roughly 10--25\,ms. Splitting would still have lost
there, by about an order of magnitude in required server slowness, so the result
does not depend on the substitute server.

\subsection{Reliability}

\begin{table}[h]
\centering
\begin{tabular}{lrrrrrr}
\toprule
Run & Requests & OK & TIMEOUT & CRC & LOST & Reboots \\
\midrule
sweep & 1{,}400 & 1{,}400 & 0 & \textbf{0} & 0 & 0 \\
\texttt{k0} & 603 & 600 & 0 & \textbf{0} & 3 & 0 \\
\texttt{k\_shallow} & 610 & 600 & 9 & \textbf{0} & 1 & 0 \\
\texttt{k\_deep} & 600 & 600 & 0 & \textbf{0} & 0 & 0 \\
\texttt{best} & 600 & 600 & 0 & \textbf{0} & 0 & 0 \\
\bottomrule
\end{tabular}
\caption{Zero corrupted deliveries across 3{,}813 requests; no node rebooted
during any run. Worst failure rate 1.6\%.}
\end{table}

Reaching this took three protocol fixes that only hardware exposed. The device
blocked for exactly the task-watchdog timeout waiting for a result, so any
request that never returned rebooted the node. A busy device silently dropped an
assignment for a different request, so one lost result cost three unacked
assignments. And the device's idle timeout outlived the orchestrator's own
deadline, leaving it deaf after the server had moved on. All three were invisible
in simulation, because in simulation nothing is ever lost.

\section{Discussion}

A single static policy is sufficient for this model on this fleet. The sweep's
argmin is \texttt{k0} for all three nodes, and the \texttt{best} policy is by
construction indistinguishable from \texttt{k0}. The orchestration machinery has
nothing to decide.

Splitting trades server work for device work, and transfer for compute. On this
fleet the exchange rate is ruinous: the head costs 121--668\,ms on the MCU to
save 2.7--3.0\,ms on the server, over a link where the entire transfer takes
7--12\,ms whichever cut produced it.

The heterogeneity is real and large. Tier B needs $4.2\times$ tier A's time for
identical work, and tier C cannot split at all. But the correct policy is
nevertheless identical on all three tiers, so the premise that motivates
per-device split selection, that different devices want different cuts, does not
hold here.

Generalising in one specific direction: split inference pays when device compute
is cheap relative to the server, and MobileNetV2-0.35 on an ESP32 is the
opposite regime. The head is 12 of 67 operations and costs 121\,ms on the best
board in the fleet; the whole model costs 3.45\,ms on one 400\,MHz x86 core. Any
argument for splitting this workload must begin by finding a server two orders
of magnitude slower, or a device two orders of magnitude faster. A much larger
model whose server cost is seconds, a far slower link, or a device with hardware
NN acceleration would each change that arithmetic, but none is a small
extrapolation from what was measured here.

\section{Threats to validity}
\label{sec:threats}

\textbf{The server is an x86 laptop, not an MCU-class edge box.} The Raspberry
Pi the design called for was unavailable. Capping the CPU to 400\,MHz brings the
tail to 3.45\,ms for the whole model, from 0.29\,ms uncapped, with a tight
distribution, but it is still far faster than a Pi. The bias favours offload;
Table~\ref{tab:sensitivity} quantifies how far it would have to move to matter,
and the answer is far enough that the conclusion survives.

\textbf{Single session, no repeats.} One run per policy. Run-to-run spread on
repeated identical runs was 10--28\% on the mean, so differences of a few percent
should not be read as real. The effects reported here are $4$--$15\times$.

\textbf{The model is a placeholder.} An ImageNet-initialised MobileNetV2-0.35
with an untrained head. Architecture, quantization and tensor shapes are those of
the real model, so all latency and byte measurements are valid; its accuracy is
meaningless and no accuracy claim is made.

\textbf{Stored images, not live capture.} All nodes read the same 50 images from
flash, which makes the comparison controlled but omits camera capture cost. The
flash read is also not quite symmetric: \texttt{k0} reads a 1.8\,KB JPEG and the
split cuts a 27{,}648\,B raw tensor, costing 17--37\,ms against 32--41\,ms,
because the cost is dominated by file opening rather than throughput. That
${\sim}10$\,ms asymmetry is reported in the \texttt{device} column and is far too
small to change any ordering.

\textbf{Loss injection was one-directional.} Injected loss reached the control
path but not the tensor fragments, so the selective-repeat path is verified in
simulation at five loss levels rather than on the radio.

\textbf{Three nodes are not radio contention.} This study says nothing about what
happens when many devices transmit simultaneously.

\section{Conclusion}

No split point beat full offload on any board, cutting deeper made things worse,
and the weakest board could not split at all. On-device inference dominates the
budget by two orders of magnitude while transfer is effectively constant across
cuts, so the parameter this project set out to tune has almost no leverage. The
negative result is quantified rather than asserted: the server would need to be
$51\times$ slower before the shallow cut broke even, roughly an order of
magnitude beyond the Pi-class server the study was designed around.

\end{document}
